% --- Archon physics formalization source begin ---
% archon:physics
% archon:covers IPhO2026Problems/problem_ipho_2026_t2_a1.lean
% archon:source-report reports/ipho_2026/problem_ipho_2026_t2_a1.source.json
% archon:problem-id ipho_2026_t2
% archon:part-id T2-A1

\chapter{Ipho Physics problem ipho\_2026\_t2\_a1}
\label{ch:IPhO2026Problems_problem_ipho_2026_t2_a1}

\paragraph{Problem source.}
\#\# Physical scenario

Parallel rays strike the inside of a half-cylindrical mirror of radius R.  For
an incident ray with transverse coordinate x, let N be its number of
reflections.  The positive threshold x\_N is the largest distance from the
optical axis for which a ray undergoes at most N reflections.  Use Figures
2c--2e for the mirror and limiting-ray geometry.

\#\# Current subquestion T2-A1

Find the general expression for the threshold x\_N in terms of R and the positive integer N.

\paragraph{Shared problem context.}
Parallel rays strike the inside of a half-cylindrical mirror of radius R.  For
an incident ray with transverse coordinate x, let N be its number of
reflections.  The positive threshold x\_N is the largest distance from the
optical axis for which a ray undergoes at most N reflections.  Use Figures
2c--2e for the mirror and limiting-ray geometry.

\paragraph{Current subquestion.}
Find the general expression for the threshold x\_N in terms of R and the positive integer N.

\paragraph{Figure/image paths.}
\begin{itemize}
\item ipho\_2026\_source/image/T2\_page-2.png
\item ipho\_2026\_source/image/T2\_page-1.png
\end{itemize}

\paragraph{Formalization target.}
create a compiling Lean file with sorry bodies at `IPhO2026Problems/problem\_ipho\_2026\_t2\_a1.lean`.
The Lean declarations must preserve the physical quantities, dimensions or dimensional roles, figure labels, governing-law hypotheses, and final relation expressed by this problem.
Use Mathlib/Physlib names found through LeanExplore where available. If a domain API is missing, introduce faithful local abstractions rather than scalar placeholder aliases.
This is an answer-blind solve-phase task. Derive a candidate result only from the problem-side evidence above; do not assume a target value, select a tolerance around a desired value, or encode a desired result into a candidate domain.
For numerical questions, define the raw end-to-end quantity and apply an explicit source-derived rounding rule. For identification questions, characterize or prove uniqueness before recording the derived candidate.

\begin{theorem}[Ipho Physics formalization target]
\label{thm:physics:ipho_2026_t2_a1:target}
The assigned autoformalize agent should translate this IPhO physics problem into Lean declarations in the covered file, with theorem and lemma proof bodies written as `by sorry`.
\end{theorem}
\begin{proof}
This is an autoformalization task, not a proof task. Produce faithful statements that can later be proved without weakening the source contract.
\end{proof}
% --- Archon physics formalization source end ---
